\documentclass[sigconf,anonymous,review]{acmart}

%% ---------------------------------------------------------------------------
%% SAC 2027, Semantic Technology track.
%% Version 2: same content and the same numbers as sac2027.tex, rewritten in
%% the authors' own writing style.
%% Double-blind: no names, no affiliation, no repository link in the source.
%% Replace \repourl with an anonymised link before submission, and with the
%% real link in the camera-ready version. It is defined once, here.
%% ---------------------------------------------------------------------------
\newcommand{\repourl}{\textsc{[anonymised repository]}}

\settopmatter{printacmref=false}
\renewcommand\footnotetextcopyrightpermission[1]{}

\usepackage{booktabs}

\begin{document}

\title{Composing Usage Policies for Models Trained on Many Sources}

\begin{abstract}
Machine learning models are increasingly trained on data held by several
organisations, each of which attaches its own conditions of use to the data it
contributes. While a number of standards exist to express such conditions for a
single dataset, and a recent ODRL profile has standardised actions such as
\emph{ai-training} and \emph{ai-augment}, none of them addresses the question
that arises as soon as more than one data owner is involved, namely what may be
done with the model that results. In this paper, we propose a composition
algebra for ODRL policies and prove that it is anti-monotone, which means that
admitting an additional contributor can never widen what the group may lawfully
do with the trained model. Since the inputs are ordinary Semantic Web artefacts,
with policies expressed as ODRL documents in JSON-LD and conditions taken from a
published ontology, composition becomes a matter of reasoning over a subsumption
hierarchy rather than of matching strings. On top of the algebra, we introduce a
measure of how similar a group's conditions are, and a selection procedure that
exploits anti-monotonicity to prune a search over candidate subsets. We then
demonstrate the approach on real data use conditions published by the European
Genome-phenome Archive, where we observe that one widely used ontology code
cannot be expressed at all, because its truth depends on the very set being
composed. Finally, we show that a group assembled the way groups are assembled
today, namely by the volume of data its members hold, is permitted to use the
model it produced only 20\% of the time, whereas the filter that avoids this
outcome requires no training at all and removes up to 98.6\% of the candidate
subsets. Code and data are available at \repourl.
\end{abstract}

\begin{CCSXML}
<ccs2012>
<concept>
<concept_id>10002951.10003260.10003282</concept_id>
<concept_desc>Information systems~Web services</concept_desc>
<concept_significance>300</concept_significance>
</concept>
<concept>
<concept_id>10002951.10003227.10003351</concept_id>
<concept_desc>Information systems~Data mining</concept_desc>
<concept_significance>300</concept_significance>
</concept>
<concept>
<concept_id>10010147.10010178</concept_id>
<concept_desc>Computing methodologies~Artificial intelligence</concept_desc>
<concept_significance>300</concept_significance>
</concept>
</ccs2012>
\end{CCSXML}

\ccsdesc[300]{Information systems~Web services}
\ccsdesc[300]{Computing methodologies~Artificial intelligence}

\keywords{ODRL, policy composition, semantic interoperability, vocabulary
alignment, linked data, usage control, data governance, Data Use Ontology,
federated learning}

\maketitle

%% ===========================================================================
\section{Introduction and Motivation}
\label{sec:intro}

Large models are increasingly adapted using private data. An institution
fine-tunes a foundation model on its own records, or builds a retrieval system
over a private collection of documents. When several institutions undertake such
an effort together, the resulting model depends on data belonging to several
owners, and each of these owners has attached its own conditions to the data it
contributed. These conditions are rarely identical, and the question of what the
group may collectively do with the trained model has no obvious answer.

To illustrate the nature of the problem, consider a simple scenario in which
three people are invited to dinner and the host is expected to respect every
guest's dietary restriction. Suppose that guest A eats only fish or vegetables,
guest B eats only meat or vegetables, and guest C eats only fish or meat. Any two
of these guests can be accommodated without difficulty, since A and B can share a
vegetable dish, A and C can share fish, and B and C can share meat. However, when
all three attend together, there is no dish that satisfies everyone, and the host
is left with nothing to serve.

Exactly the same structure arises in collaborative model training. Suppose that
three institutions declare the purposes for which their data may be used, such
that contributor A permits cancer research or method development, contributor B
permits neurological research or method development, and contributor C permits
cancer or neurological research. Any two of these contributors share a purpose,
yet all three of them share none. A consortium that assessed each prospective
member separately against its own research question would therefore conclude that
every member is acceptable, and would be mistaken.

This small scenario already exhibits the three properties that the remainder of
this paper formalises. First, conditions intersect, and consequently the set of
actions available to the group can only shrink as further members are admitted.
Second, and as a direct consequence, admitting a member can never introduce a
permission that was not already present, which is the anti-monotonicity we
establish in Section~\ref{sec:approach}. Third, compatibility is not transitive,
which means that verifying members one at a time proves nothing, since the
problem becomes visible only in the set of three. For this reason, the selection
procedure we propose searches over subsets rather than filtering candidates
individually.

A number of standards for expressing conditions of use already exist and are
actively used in practice. SPDX~\cite{spdx} describes software licences and
supports expressions that combine them, the GA4GH Data Use Ontology
(DUO)~\cite{duo} describes conditions attached to biomedical datasets and is
published alongside real archives, and the W3C Open Digital Rights Language
(ODRL)~\cite{odrl-model,odrl-vocab} provides a general model of permissions,
prohibitions, duties and constraints. More recently, in September 2026, the ODRL
community published a draft AI profile~\cite{odrl-ai} that introduces actions
specific to AI workflows, namely \emph{data-mining}, \emph{ai-training},
\emph{gen-ai}, \emph{ai-augment}, \emph{search} and \emph{index}. All of these
efforts describe what may be done to a dataset, yet none of them specifies what
happens when several policies have to hold simultaneously, and none of them
describes the artefact that training produces. This is precisely the gap that our
work addresses, and the question is far from hypothetical, since the new ODRL
profile defines \emph{ai-training} for fine-tuning and \emph{ai-augment} for
retrieval and grounding, which are exactly the two ways in which foundation
models are adapted with private data today.

We approach this as a problem for semantic technology, because that is what it
is. The policies involved are linked data, expressed as ODRL documents in JSON-LD
whose actions, operands and values are terms identified by IRIs and drawn from
published vocabularies. The conditions attached to the data are classes of an
ontology with stable identifiers, and the purposes these classes name form a
hierarchy rather than a flat list. What we contribute is a semantics for
combining such documents, together with a mapping between two independently
designed vocabularies, ODRL and DUO, in which the fidelity of every
correspondence is recorded. As Section~\ref{sec:approach} demonstrates, treating
the vocabulary properly is not merely a presentational choice, since composition
that compares terms as strings rather than reasoning over their subsumption
hierarchy reports conflicts that do not exist, and would have altered the results
we report.

It is worth delimiting the scope of our contribution at the outset. Selecting a
group of contributors consists of two distinct parts, namely determining which
groups are permitted and determining which of the permitted groups produces the
best model. This paper addresses the first of these. We assume that a utility
measure is supplied by the application, and we make no claim as to which
permitted group trains best. As Section~\ref{sec:poc} shows, solving the first
part is in fact what makes the second part affordable.

Against this background, this paper makes two contributions.

\begin{enumerate}
\item First, we propose an approach to composing usage policies which is
generic, in that it depends neither on a particular application domain nor on a
particular vocabulary. It consists of the following three elements:
  \begin{enumerate}
  \item a composition algebra for a fragment of ODRL, together with a proof that
  capability is anti-monotone in the contributor set, so that admitting a
  contributor can only narrow what the group may do;
  \item a measure of how similar a group's conditions are, which separates
  genuine shared practice from agreement arising by coincidence, together with a
  selection procedure that exploits anti-monotonicity in order to prune a search
  over candidate subsets; and
  \item a model bundle that carries its composed agreement, in which
  verification recomputes the agreement from the recorded declarations rather
  than comparing a hash, and in which composition is recursive, so that reuse of
  the model can only narrow the rights attached to it.
  \end{enumerate}
\item Second, we provide a proof of concept in which the approach is applied to
real data use conditions published by a public archive. It contributes:
  \begin{enumerate}
  \item a mapping of GA4GH DUO as published by the European Genome-phenome
  Archive, covering all 28 codes observed there, including one that cannot be
  expressed at all, together with an explanation of why this is so; and
  \item a quantification of the cost of ignoring policy, showing that a group
  assembled the way groups are assembled today is permitted to use its own model
  only 20\% of the time, which means that most of the computation spent on
  training it is wasted.
  \end{enumerate}
\end{enumerate}

The rest of this paper is organised as follows. Section~\ref{sec:related}
positions our work with respect to existing standards and research efforts, and
explains what none of them currently provides. Section~\ref{sec:approach} then
presents the proposed approach, which is generic and depends neither on a
particular domain nor on a particular vocabulary. Section~\ref{sec:poc} applies
the approach to real data use conditions published by a public archive and
reports what it finds and what it saves. Section~\ref{sec:discussion} discusses
the results together with the assumptions and limitations of the work, and
Section~\ref{sec:conclusion} concludes the paper with final remarks and an
outlook for future work.

%% ===========================================================================
\section{Background and Related Work}
\label{sec:related}

A considerable body of work describes conditions of use, and some of it is
deployed at scale. None of it, however, composes conditions belonging to several
owners, and none of it describes the model that training produces. In this
section we review what exists, so that the gap we address becomes visible rather
than merely asserted.

\subsection{ODRL and the Data Use Ontology}

ODRL is a W3C Recommendation and forms part of the Semantic Web stack, in the
sense that a policy is an RDF graph, usually serialised as JSON-LD, in which
every action, operand and value is a term with its own IRI. A policy is expressed
as a set of rules, where a \emph{permission} allows an action, a
\emph{prohibition} forbids one, and a \emph{duty} constitutes an obligation that
must be discharged. Rules may carry \emph{constraints}, which restrict the
circumstances under which a rule applies, for instance to a particular purpose,
region or date. In our implementation, policies are read by a parser for the
JSON-LD serialisation, which has to accommodate the forms encountered in real
documents, where terms appear bare, prefixed or as full IRIs, and where a single
object may stand wherever a list is permitted.

We employ a fragment of ODRL comprising five actions relevant to model building,
namely \emph{use}, \emph{aggregate}, \emph{derive}, \emph{distribute} and
\emph{transfer}, of which the first three are the most significant. A model
trained on several contributions is an aggregation of them, and subsequently
fine-tuning that model constitutes a derivation. ODRL provides both of these in
its core vocabulary, which is the main reason it fits this problem better than a
general-purpose access control language. The draft AI profile~\cite{odrl-ai}
additionally introduces actions arranged in a hierarchy, in which \emph{gen-ai}
is a kind of \emph{ai-training}, which is in turn a kind of \emph{data-mining}.
Our fragment accommodates these actions directly, because the algebra treats
actions through a subsumption relation rather than as flat labels.

DUO~\cite{duo}, in turn, is an ontology of conditions attached to biomedical
data, published with stable term identifiers and adopted by the European
Genome-phenome Archive (EGA)~\cite{ega} among others. A dataset is annotated with
codes such as DUO:0000006 (HMB, health or medical or biomedical research only),
NCU (no commercial use), or IRB (ethics approval required). Relating these terms
to ODRL constitutes a vocabulary alignment problem, and one in which the two
vocabularies were designed for different purposes, with the consequence that some
correspondences are exact while others are interpretations. We record which is
which.

\subsection{Describing Conditions on a Single Asset}

DUO has a companion in ADA-M~\cite{adam}, which was developed to make data use
conditions automatable and is therefore the closest prior work to our own.
Both, however, are built around matching a single request against a single
dataset, and neither provides a notion of combining several datasets in order to
determine what the joint rights are.

A second group of formats documents datasets and models for human readers.
Datasheets for Datasets~\cite{datasheets} and Model Cards~\cite{modelcards}
require authors to record provenance, intended use and known limitations, while
Croissant~\cite{croissant} supplies machine-readable metadata for datasets,
including licence fields, and is used by large dataset repositories. The W3C Data
Privacy Vocabulary~\cite{dpv} provides controlled terms for purposes, and the
ODRL AI profile adopts it for exactly this role. These efforts are valuable, yet
they remain descriptive, in that they record what went into a system without
defining any operation that derives what comes out of it.

That this matters in practice is not a matter of opinion. A recent audit of
widely used machine learning datasets found licence information on aggregator
sites to be missing or incorrect at a high rate~\cite{provenance}, and our own
measurement points in the same direction, since of the 4\,817 policy records
available in EGA only 21.4\% carry machine-readable conditions at all, as
discussed in Section~\ref{sec:corpus}.

\subsection{Licences that Travel with Derivatives}

Software licence composition represents the mature analogue of our problem.
SPDX~\cite{spdx}, now an ISO standard, expresses licences and combines them using
operators, and copyleft propagation is the very phenomenon we study, in that
combining incompatible terms yields an artefact that may not be distributed.
Bills of materials subsequently turned this into established practice, and
CycloneDX~\cite{cyclonedx} extends the idea to models and datasets through a
machine learning bill of materials. These approaches record the ingredients,
however, and do not derive what may be done with the result.

The OpenRAIL family~\cite{openrail} attaches use restrictions that downstream
users are required to pass on, and is deployed at scale on widely used models,
although the mechanism is contractual and intended for human readers. Nothing in
it computes the outcome of combining two such licences carrying different
restriction lists. C2PA~\cite{c2pa}, finally, signs provenance assertions so that
the history of a piece of content can be verified cryptographically, which makes
it the natural prior art for the signing gap we describe in
Section~\ref{sec:bundle}, although it addresses integrity rather than the meaning
of combined terms.

\subsection{Usage Control and Data Spaces}

UCON~\cite{ucon} and the subsequent usage control literature study obligations
that persist after access has been granted, with decisions that are continuous
rather than taken once. This body of theory underpins our treatment of duties. It
is also the setting in which ODRL is most widely used commercially, since
industrial data spaces~\cite{ids} apply it to data shared between companies,
using terms such as permitted purpose, retention period and region. The emphasis
in that work, however, lies on enforcing terms attached to a single asset as it
moves between two parties, rather than on combining the terms of many parties
over a shared artefact.

\subsection{Opt-Out Vocabularies for AI}

A recent group of proposals allows publishers to state whether their content may
be used for AI training, among them the IETF AI preferences
vocabulary~\cite{aipref} and several industry formats. The draft ODRL AI
profile~\cite{odrl-ai} reconciles these proposals, observing that they are
divergent and fragmented. Regulation is pushing in the same direction, in that
the EU AI Act~\cite{aiact} obliges providers of general-purpose models to
summarise their training content, while the Digital Single Market
directive~\cite{dsm} permits rights holders to reserve text and data mining
rights in machine-readable form. All of this work, however, concerns a single
publisher and a single asset, and therefore addresses only the input side of the
problem.

\subsection{What None of This Provides}

The pattern emerging from this review is consistent. Software licence composition
is a solved problem, yet it concerns software. Data use conditions are machine
readable, yet they are matched one dataset at a time. AI opt-out vocabularies
standardise the action, yet they stop at the asset. Documentation formats record
what went in, yet they define no operation over it. Usage control enforces the
terms of a single agreement between two parties. Consequently, two things are
missing from all of them, namely a semantics for combining the conditions of
several owners, and a vocabulary for the artefact that training produces. Both
become necessary as soon as a model is built from more than one governed source,
which is now the normal case.

%% ===========================================================================
\section{Proposed Approach}
\label{sec:approach}

Our proposed approach is generic. It assumes only that contributors state their
conditions as ODRL policies, and that the terms used in those conditions are
drawn from a vocabulary equipped with a subsumption hierarchy. Nothing in this
section depends on the application domain, on biomedical data, or on any
particular vocabulary, and Section~\ref{sec:poc} subsequently instantiates it.

\subsection{Composition Rules}

Let $S$ denote a set of contributors, each of which holds a policy $p_i$. The
composed agreement $A(S)$ is constructed element by element, following the rules
summarised in Table~\ref{tab:rules}. Permissions intersect, because an action
requires the grant of every contributor; prohibitions are unioned, because a
single prohibition binds the entire group; duties are likewise unioned and merged
per action, so that the strictest obligation applies; constraints are conjoined
through a set meet or an interval intersection; and conflicts are resolved by
letting denial override.

\begin{table}[h]
\caption{Composition rules. Every rule narrows.}
\label{tab:rules}
\begin{tabular}{ll}
\toprule
Element & Rule \\
\midrule
Permissions & intersect: an action needs every contributor's grant \\
Prohibitions & union: one prohibition binds the whole group \\
Duties & union, merged per action: the strictest applies \\
Constraints & conjoin: set meet, or interval intersection \\
Conflict & deny overrides \\
\bottomrule
\end{tabular}
\end{table}

\subsection{Anti-Monotonicity}

Writing $\mathrm{Cap}(S)$ for the set of actions the group may perform under
$A(S)$, the central property of the algebra is that, for contributor sets
$S \subseteq S'$,
\[
\mathrm{Cap}(S') \subseteq \mathrm{Cap}(S).
\]

The argument follows directly from the composition rules. Since permissions
intersect, an action permitted by $A(S')$ must be permitted by every member of
$S'$, and therefore by every member of $S$. Prohibitions and duties only
accumulate, and constraints only narrow. It follows that admitting a contributor
can never widen what the group may do with the resulting model, however
permissive that contributor may be, and that it can only add data. Conversely,
removing a contributor can only widen the group's capability. This is the dinner
scenario of Section~\ref{sec:intro} stated in general terms, in that admitting
a further guest never adds a dish to the menu. In addition to the proof, we verify
the property through randomised property-based testing over generated policies,
since it is the implementation rather than the proof that our experiments rely
on.

\subsection{The Role of Subsumption}

Constraints must be compared through a subsumption relation rather than as
strings. A contributor permitting use within the EU and a contributor permitting
use within Germany are compatible, and the meet of their constraints is Germany,
whereas under a flat vocabulary the two would be read as a conflict and the
capability would be lost. This is not a matter of presentation. Were the
vocabulary flat, part of the capability loss measured in Section~\ref{sec:poc}
would be an artefact of our own modelling rather than a property of the policies
themselves. The same consideration applies to the action hierarchy of the AI
profile, in which a contributor permitting \emph{data-mining} and a contributor
permitting only \emph{ai-training} must compose to \emph{ai-training}.

\subsection{A Similarity Measure for a Group}
\label{sec:similarity}

The extent to which a group loses capability depends on how alike its members'
conditions are, and we therefore make this a measurable quantity rather than an
assumption. We describe similarity by means of a single parameter $c$ taking
values between 0 and 1, which answers the question of how likely two institutions
belonging to the same group are to have written the same conditions. At $c=0$,
every institution decides independently and two of them coincide only by
accident, whereas at $c=1$ every institution follows a common template and any
two of them always coincide. Real groups lie somewhere between these extremes.

This parameter is what governs how rapidly capability falls. If all contributors
hold identical conditions, combining them costs nothing, since intersecting a set
with itself returns that same set, whereas if their conditions differ, every
additional contributor removes something further. It is therefore $c$, rather
than the number of contributors, that determines the outcome.

The parameter cannot be observed directly, however. What can be observed is how
often two policies happen to be identical, and this is not the same quantity,
because two institutions may coincide purely by chance, given that some condition
sets are common and two institutions deciding independently will occasionally
write the same thing. The observed rate consequently mixes two causes, namely
genuine shared practice and coincidence, and the model we propose separates them.
Writing $p$ for the coincidence rate, that is, the frequency with which two
independent draws from the observed distribution of condition sets produce the
same set, the probability that the condition sets of two contributors are
identical is given by
\[
\mathrm{agreement} = c^2 + (1-c)^2 p + 2c(1-c)w,
\]
where $w$ denotes the weight of the shared set. Since the shared set is itself
drawn from the same distribution, we have $\mathbb{E}[w]=p$, and the expression
reduces to
\[
\mathrm{agreement} = p + c^2 (1-p),
\qquad
c = \sqrt{\frac{\mathrm{agreement}-p}{1-p}}.
\]

Both $p$ and the observed agreement are counted from a corpus of policies, so
that $c$ follows from the data. Agreement at or below $p$ yields $c=0$, since the
model is able to add similarity but not to subtract it, and it therefore reports
an absence of shared practice rather than a negative value.

Applying the measure requires a notion of which policies belong to the same
group, and any observable grouping will inevitably be imperfect. For this reason
we do not report a single value. Instead, we compute $c$ for the tightest
grouping the data offers, which yields an upper bound, and for unrelated
policies, which yields a lower bound, and we report the interval between the two.
Section~\ref{sec:measured} carries this out for a real archive.

\subsection{Selecting a Group}
\label{sec:selection}

Given that composing an arbitrary group usually yields a model that nobody may
use, the practical question becomes how to find a group that does not suffer from
this problem. As the scenario in Section~\ref{sec:intro} illustrated, this is a
search rather than a filter, since compatibility is not transitive, feasibility is
a property of the set as a whole, and candidates therefore cannot be checked
individually against a reference policy.

Anti-monotonicity is what makes the search tractable. Because capability only
narrows as the set grows, feasibility is downward closed, in the sense that every
subset of a feasible set is itself feasible. A contributor that fails on its own
can consequently be discarded outright, and a set of size $k+1$ need only be
considered once all of its subsets of size $k$ are known to be feasible. This is
precisely the candidate generation step of the Apriori algorithm~\cite{apriori},
and the pruning is sound because the property it relies upon is exact rather than
estimated.

The filter itself costs nothing, in that it requires no training and no data
leaves any site. It can therefore be applied first, so that whatever method is
subsequently used to rank the remaining groups needs to consider only those
groups that could actually be deployed. Training a group that may not use its own
model is wasted computation, and the ordering is consequently not a matter of
convenience but the only ordering that makes sense. We do not ourselves rank the
permitted groups, which is not to say that they are equally good, but rather that
choosing among them requires a measure of expected model quality that this paper
neither provides nor assumes. What the filter guarantees is that whatever measure
is applied later will be applied to groups that could lawfully be deployed, and
we return to this point in Section~\ref{sec:conclusion}.

\subsection{Models that Carry their Policy}
\label{sec:bundle}

A trained model is a derivative of every contributor's data, and the composed
agreement states what may lawfully be done with it. Should the two be separated,
the terms are lost the first time the model file is copied. We therefore export a
bundle consisting of the weights together with a manifest that records the
agreement, every contributor declaration from which it was composed, and whatever
obligations remain outstanding.

Verifying such a bundle does not amount to comparing a digest against itself.
Instead, the agreement is recomposed from the declarations recorded in the
manifest and checked against what the manifest claims, after which the weights
are confirmed to still hash to the recorded value. Widening a purpose therefore
requires forging every contributor's declaration rather than editing a single
line. Composition is moreover recursive, since an agreement is itself a policy,
with the consequence that a bundle's agreement composes back in whenever the
model is reused as an input. Because capability can only narrow, a restricted
model cannot be turned into an unrestricted one by training something else upon
it, and it is this property that makes a policy-carrying model more than mere
documentation, and that none of the formats reviewed in
Section~\ref{sec:related} provides.

The manifest is not signed, and we state this explicitly rather than leaving it
to be discovered. A forger who rewrites the declarations and the agreement
consistently will produce a bundle that verifies, so that detection covers
accidental corruption and one-sided edits, while closing the remaining gap would
require signatures over the digest together with a key distribution mechanism,
for which C2PA~\cite{c2pa} offers the obvious model.

%% ===========================================================================
\section{Proof of Concept}
\label{sec:poc}

In this section we apply the proposed approach to real conditions published by a
real archive, in order to show both what it finds and what it saves.

\subsection{The Corpus}
\label{sec:corpus}

We take the distribution of conditions from EGA policy records, which are
publicly available and require no registration. Of the 4\,817 policy records
examined, 1\,031 carry structured DUO codes, which corresponds to an adoption
rate of 21.4\%. This figure is itself a result about current practice rather than
an inconvenience, since the majority of data use conditions in this archive are
expressed as free text or as an external link, and are therefore not machine
readable at all. We rely on these records rather than on invented policies
because the distribution of conditions encountered in the wild is itself part of
the result. Had we invented a distribution, the main measurement would have been
unanswerable, in that a reader could reasonably ask why capability falls as it
does, and the only honest answer would be that we chose it.

Mapping DUO into the ODRL fragment requires every code to be classified, and we
classify all 28 codes observed in the archive, as summarised in
Table~\ref{tab:duo}. Three of these are deprecated, having been retired from the
ontology while remaining present in live annotations, which means that any tool
encountering real governance metadata has to decide what to do with obsolete
terms. Every mapping furthermore carries a fidelity
marker, either \emph{exact} or \emph{interpretive}, and the distribution of
these markers is itself informative: all of the restrictions map exactly, whereas
the externally adjudicated codes and the duties are without exception
interpretations, which is where the two vocabularies were designed for different
ends. The complete mapping, code by code, is available in the repository.

\begin{table}[h]
\caption{Classification of the 28 DUO codes observed in EGA.}
\label{tab:duo}
\begin{tabular}{llrrrr}
\toprule
 & & & \multicolumn{3}{c}{Fidelity} \\
\cmidrule(l){4-6}
Class & Meaning & $n$ & exact & interp. & n/a \\
\midrule
permission & sets a purpose & 8 & 7 & 1 & --- \\
restriction & constrains or forbids & 8 & 8 & --- & --- \\
externally adjudicated & decided out of band & 4 & --- & 4 & --- \\
deprecated & obsolete, still in use & 3 & --- & 2 & 1 \\
duty & an obligation & 2 & --- & 2 & --- \\
structural & ontology category & 2 & --- & --- & 2 \\
out of fragment & inexpressible & 1 & --- & --- & 1 \\
\midrule
total & & 28 & 15 & 9 & 4 \\
\bottomrule
\end{tabular}
\end{table}

One code, however, cannot be expressed at all. DUO:0000020 (COL, collaboration
required) states that the data provider must be a participant in the work, so
that whether it is satisfied depends on the very contributor set being composed.
This is self-referential, and it is the one construct our fragment excludes by
design. We report it rather than conceal it, since a widely used ontology
containing a condition that a compositional treatment cannot accept without a
fixed point is a finding in its own right, and anyone building upon DUO will
encounter the same difficulty.

\subsection{Measured Similarity}
\label{sec:measured}

EGA policy records name the Data Access Committee that issues them, and the
1\,031 annotated policies originate from 899 distinct committees. The committee
therefore provides the grouping that Section~\ref{sec:similarity} requires. For
this corpus the coincidence rate is $p = 0.0131$, which means that two unrelated
institutions agree approximately 1.3\% of the time for no reason other than
chance, and any agreement above this rate constitutes evidence of shared
practice.

\begin{table}[h]
\caption{Similarity measured from EGA policy records.}
\label{tab:cohesion}
\begin{tabular}{lrl}
\toprule
Quantity & Value & Implies \\
\midrule
chance coincidence $p$ & 0.0131 & --- \\
agreement, different committees & 0.0118 & $c=0.000$ \\
agreement, same committee & 0.474 & $c=0.683$ \\
\bottomrule
\end{tabular}
\end{table}

The two rows of Table~\ref{tab:cohesion} read very differently. Policies issued
by two different committees are identical 1.18\% of the time, which lies slightly
below the 1.31\% that chance alone would produce, so that there is no evidence of
shared practice between institutions and $c$ is zero. Policies issued by the same
committee, by contrast, are identical 47.4\% of the time, which is far above
chance and yields $c = 0.683$.

A committee, however, is not a consortium. It is a single institution issuing
several policies, usually written by the same people under the same internal
rules, and it is therefore the most similar grouping that can be found anywhere
in the archive. A real consortium consists of different institutions and must
consequently be less similar than this. The same-committee figure accordingly
serves as an upper bound and the across-committee figure as a lower bound, and we
report the interval rather than selecting a point within it.

\subsection{Capability Loss}

We draw contributor policies from the observed distribution and compose subsets of
increasing size, measuring the proportion of groups that retain any usable
permission. Table~\ref{tab:erosion} reports this at both ends of the measured
interval, and additionally for a corpus that assumes nothing, in which condition
sets are drawn whole from the committees that published them so as to preserve
whatever structure the archive actually exhibits.

\begin{table}[h]
\caption{Capability lost at ten contributors.}
\label{tab:erosion}
\begin{tabular}{lrr}
\toprule
Condition & Loss at 10 & Retains any \\
\midrule
measured ceiling ($c=0.683$) & 16\% & 84\% \\
drawn from real committees & 52\% & 48\% \\
measured floor ($c=0.000$) & 53\% & 47\% \\
\midrule
assumed baseline, for reference & 64\% & 36\% \\
\bottomrule
\end{tabular}
\end{table}

These figures matter because the alternative is to assume. Across seven assumed
similarity models, the loss at ten contributors ranges from 7\% to 77\%, which
covers most of the possible range, whereas measuring narrows that range from 69
percentage points to 37. The loss is substantial under every condition, so that
the direction of the result does not depend on the modelling choice, although the
magnitude does and must therefore always be quoted together with the condition
that produced it. The upper bound, moreover, rests on the 57 committees that
publish more than one machine-readable policy, which is a small number and is
itself a finding, since alongside the adoption rate of 21.4\% it indicates that
structured data use policy is not yet common practice.

\subsection{Selection and the Computation it Saves}
\label{sec:saving}

Evaluating selection requires a baseline, and the comparison is only meaningful
if that baseline reflects how groups are assembled in practice. Our policy-blind
baseline is accordingly not a random group. Contributors are ranked by the volume
of data they bring, and the best $N$ of them are taken, with $N$ set to the size
returned by the policy-aware search, so that both groups contain the same number
of members and differ only in whether policy was consulted. This is deliberately
the strongest form of the baseline, since ranking contributors by the data they
hold is standard practice, and a weaker baseline would make the comparison easier
than it ought to be.

We run the search of Section~\ref{sec:selection} over a pool of 16 contributors,
averaged over 5 generated corpora, under three requests of increasing strictness.
Table~\ref{tab:selection} reports two quantities, of which \emph{space removed}
is the proportion of the $2^{16}-1$ candidate subsets that the pruning never had
to examine, while \emph{blind usable} is the frequency with which the
policy-blind group turns out to be permitted to use the model it produced.

\begin{table}[h]
\caption{Selection results by request.}
\label{tab:selection}
\begin{tabular}{lrr}
\toprule
Request & Space removed & Blind usable \\
\midrule
actions only & 60.9\% & 20\% \\
purpose = cancer research & 83.7\% & 0\% \\
purpose = research & 98.6\% & 0\% \\
\bottomrule
\end{tabular}
\end{table}

The second column is best read as a waste rate. A group chosen by data volume
alone is permitted to use its own model 20\% of the time under the loosest
request, and never once the request names a purpose, which means that four out of
five training runs are wasted in the first case and all of them in the remaining
two. This waste is furthermore not detected by training itself, since the model
converges, reports an accuracy, and remains unusable for a reason that no metric
reveals. The filter that avoids this outcome, by contrast, costs nothing, in that
it requires no training, moves no data, and is exact rather than estimated, so
that a group it rejects is genuinely not permitted. It additionally shrinks the
problem for whatever follows, since between 60.9\% and 98.6\% of the $2^{16}-1$
candidate subsets are removed without ever being examined.

Finally, it is worth observing that the pruning performs best precisely where
selection is hard, and degenerates towards full enumeration where the requirement
is so permissive that almost any group would have satisfied it. This is a
property of the problem rather than a weakness of the algorithm, since when the
constraint is weak the family of feasible sets genuinely is that large.

%% ===========================================================================
\section{Discussion}
\label{sec:discussion}

\subsection{Results}

Our results support three observations, listed here in decreasing order of how
firmly they are established. First, composition loses a substantial amount of
capability, and how much it loses depends on similarity rather than on group
size, with the loss at ten contributors ranging from 16\% to 53\% across the
measured interval. This is a property of the declared terms and not of any model.
Second, choosing a group without consulting policy wastes most of the computation
spent upon it, because the resulting model usually may not be used, as
Section~\ref{sec:saving} shows, whereas the filter that prevents this is free and
additionally removes most of the search space. Third, similarity can be measured
rather than assumed, which narrows the range of plausible answers from 69
percentage points to 37, although it cannot be reduced to a single value given
the data available, and we therefore report an interval instead.

\subsection{Assumptions}

Our approach assumes that contributors declare their conditions truthfully, and
this assumption is reasonable for a reason worth making explicit. Overstating a
restriction only reduces what the declaring party may itself do, so that there is
no incentive to do so. The same reasoning does not extend to other properties a
contributor might declare, which is precisely why we do not ask contributors to
declare any. We further assume that conditions are expressed in a vocabulary
equipped with a subsumption hierarchy, since without one composition reports
conflicts that do not exist, as discussed in Section~\ref{sec:approach}. Finally,
we assume that the data held by contributors is of comparable quality, because
the method observes conditions rather than data, with the consequence that a
contributor combining permissive terms with poor data remains invisible to it and
would have to be excluded by other means.

\subsection{Limitations}

Several limitations of the present work should be acknowledged. Bundles are not
signed, as discussed in Section~\ref{sec:bundle}. The similarity upper bound
rests on 57 committees, and although it is reported with a bootstrap interval,
only the lower bound constitutes a hard floor rather than an estimate. Model
quality is not addressed at all, in that we make no claim as to which permitted
group produces the better model, and while our evaluation of capability never
touches a metric, which is what makes it stable, it also means that this paper
says nothing about accuracy. Lastly, most of the work is not specific to
federated learning, since the algebra applies to any derivative built from
several governed sources, whether a trusted research environment, a shared
analysis platform, or a model fine-tuned on pooled records. Federated
learning~\cite{fedavg} is simply where we implemented it, because the policy
check can there be run inside the system before training begins, and not because
the result depends upon it. Centralising the data would not avoid the problem
either, since the model remains a derivative of every source and the terms of
every source continue to bind it.

%% ===========================================================================
\section{Conclusion and Future Work}
\label{sec:conclusion}

Conditions attached to data compose, and they compose in a single direction. In
this paper we proposed an algebra for a fragment of ODRL, proved that capability
is anti-monotone in the contributor set, and built two further components upon
it, namely a measure of how similar a group's conditions are and a selection
procedure that exploits anti-monotonicity to prune a search over candidate
subsets. Applied to a public archive of real data use conditions, the approach
shows that a group assembled the way groups are assembled today is usually not
permitted to use the model it produced, and that the filter preventing this
outcome requires no training whatsoever.

Ranking the permitted groups represents the obvious next step, and the one we did
not take. Several directions are available for this purpose. Contribution
measures such as Data Shapley~\cite{datashapley}, influence functions or
leave-one-out estimate what each contributor adds, albeit at a cost that grows
rapidly with group size. Client selection methods originating in federated
learning employ cheaper signals, with Oort~\cite{oort} ranking participants by a
statistical utility derived from their local loss. Short exploratory runs
constitute a third option, in which every permitted group is trained for a few
rounds, ranked, and only the best are continued.

One asymmetry makes this harder than it may appear. Permissions can be taken from
a declaration, because claiming more restrictions than one actually has only
shrinks what one may do, whereas data quality cannot, because the incentive runs
in the opposite direction. Quality therefore has to be measured, and measuring it
is the expensive problem that the filter renders affordable rather than solves.
Signing the manifest constitutes the other open item, and the one that would turn
policy-carrying models from a description into a mechanism.

\subsection*{Availability}

The implementation, the policy corpus, the committee metadata and every result
table reported in this paper are available at \repourl. All tables are produced
from committed result files, so that every number reported here can be traced
back to the run that produced it.

\bibliographystyle{ACM-Reference-Format}
\bibliography{refs}

\end{document}
